\exercisegroup

\begin{enumerate}
\def\labelenumi{\arabic{enumi}.}
\item
  Let E be an ordered set in which there exists at least one pair of distinct comparable elements. Show that, if R \(\{x, y\}\) denotes the relation `` \(x \in E\) and \(y \in E\) and \(x < y\) '', then R satisfies the first two conditions of no. 1 but not the third.
\item
  \begin{enumerate}
  \def\labelenumii{(\alph{enumii})}
  \tightlist
  \item
    Let E be a preordered set and let S\{x, y\} be an equivalence relation on E. Let R\{X, Y\} denote the relation ``X ∈ E/S and Y ∈ E/S and for each x ∈ X there exists y ∈ Y such that x ≤ y''. Show that R is a preorder relation on E/S, called the quotient by S of the relation x ≤ y. The quotient set E/S, endowed with this preorder relation, is called (by abuse of language; cf.~Chapter IV, § 2, no. 6) the quotient by S of the preordered set E.
  \end{enumerate}
\end{enumerate}

\begin{enumerate}
\def\labelenumi{(\alph{enumi})}
\setcounter{enumi}{1}
\tightlist
\item
  Let \(\varphi\) be the canonical mapping of E onto E/S. Show that if \(g\) is a mapping of the preordered quotient set E/S into a preordered set F such that \(g \circ \varphi\) is an increasing mapping, then \(g\) is an increasing mapping. The mapping \(\varphi\) is increasing if and only if S satisfies the following condition:
\end{enumerate}

\begin{enumerate}
\def\labelenumi{(\Alph{enumi})}
\setcounter{enumi}{2}
\tightlist
\item
  The relations \(x \leqslant y\) and \(x \equiv x' \pmod{S}\) in \(E\) imply that there exists \(y' \in E\) such that \(y \equiv y' \pmod{S}\) and \(x' \leqslant y'\) .
\end{enumerate}

If this condition is satisfied, the equivalence relation S is said to be weakly compatible (in x and y) with the preorder relation \(x \leqslant y\) . Every equivalence relation S which is compatible (in x) with the preorder relation \(x \leqslant y\) (Chapter II, §6, no. 3) is a fortiori weakly compatible (in x and y) with this relation.

\begin{enumerate}
\def\labelenumi{(\alph{enumi})}
\setcounter{enumi}{2}
\item
  Let \(E_{1}\) and \(E_{2}\) be two preordered sets. Show that if \(S_{1}\) is the equivalence relation \(pr_{1}z = pr_{1}z'\) on \(E_{1} \times E_{2}\) , then \(S_{1}\) is weakly compatible in \(z\) and \(t\) with the product preorder relation \(z \leqslant t\) on \(\mathbf{E}_1 \times \mathbf{E}_2\) (but is not usually compatible with this relation in \(z\) or \(t\) separately); moreover, if \(\varphi_1\) is the canonical mapping of \(\mathbf{E}_1 \times \mathbf{E}_2\) onto \((\mathbf{E}_1 \times \mathbf{E}_2)/\mathbf{S}_1\) , and if \(\mathrm{pr}_1 = f_1 \circ \varphi_1\) is the canonical decomposition of \(\mathrm{pr}_1\) with respect to the equivalence relation \(\mathbf{S}_1\) , then \(f_1\) is an isomorphism of \((\mathbf{E}_1 \times \mathbf{E}_2)/\mathbf{S}_1\) onto \(\mathbf{E}_1\) .
\item
  With the hypothesis of (a), suppose that E is an ordered set and that the following condition is satisfied :
\end{enumerate}

(C') The relations \(x \leqslant y \leqslant z\) and \(x \equiv z \pmod{S}\) in E imply \(x \equiv y \pmod{S}\) .

Show that \(\mathbf{R}\{\mathbf{X},\mathbf{Y}\}\) is then an order relation between \(\mathbf{X}\) and \(\mathbf{Y}\) on \(\mathbf{E} / \mathbf{S}\) .

\begin{enumerate}
\def\labelenumi{(\alph{enumi})}
\setcounter{enumi}{4}
\item
  Give an example of a totally ordered set E with four elements and an equivalence relation S on E such that neither of the conditions (C) and (C') is satisfied, but such that E/S is an ordered set.
\item
  Let E be an ordered set, let f be an increasing mapping of E into an ordered set F, and let \(S\{x,y\}\) be the equivalence relation \(f(x)=f(y)\) on E. Then the condition \((\mathbf{C}')\) is satisfied. Moreover, the condition (C) is satisfied if and only if the relations \(x\leqslant y\) and \(f(x)=f(x')\) imply that there exists \(y'\in E\) such that \(x'\leqslant y'\) and \(f(y)=f(y')\) . Let \(f=g\circ\varphi\) be the canonical decomposition of f.~Then g is an isomorphism of E/S onto \(f(\mathbf{E})\) if and only if this condition is satisfied and, in addition, the relation \(f(x)\leqslant f(y)\) implies that there exist \(x',y'\) such that \(f(x)=f(x')\) , \(f(y)=f(y')\) , and \(x'\leqslant y'\) .
\end{enumerate}

\begin{enumerate}
\def\labelenumi{\arabic{enumi}.}
\setcounter{enumi}{2}
\tightlist
\item
  Let I be an ordered set and let \((\mathbf{E}_i)_{i\in \mathbf{I}}\) be a family of non-empty ordered sets indexed by I.
\end{enumerate}

\begin{enumerate}
\def\labelenumi{(\alph{enumi})}
\item
  Let F be the sum (Chapter II, § 4, no. 8) of the family \((\mathbf{E}_t)_{t \in I}\) ; for each \(x \in F\) , let \(\lambda(x)\) be the index \(t\) such that \(x \in \mathbf{E}_t\) ; and let G be the graph consisting of all pairs \((x, y) \in F \times F\) such that either \(\lambda(x) < \lambda(y)\) or else \(\lambda(x) = \lambda(y)\) and \(x \leqslant y\) in \(\mathbf{E}_{\lambda(x)}\) . Show that G is the graph of an ordering on F. The set F endowed with this ordering is called the ordinal sum of the family \((\mathbf{E}_t)_{t \in I}\) (relative to the ordering on I) and is denoted by \(\sum_{t \in I} \mathbf{E}_t\) . Show that the equivalence relation corresponding to the partition \((\mathbf{E}_t)_{t \in I}\) of F satisfies conditions (C) and (C') of Exercise 2, and that the quotient ordered set (Exercise 2) is canonically isomorphic to I.
\item
  If the set I is the ordinal sum of a family \((J_{\lambda})_{\lambda\in L}\) of ordered sets, where L is an ordered set, show that the ordered set \(\sum_{i\in I}E_{i}\) is canonically isomorphic to the ordinal sum \(\sum_{\lambda\in L}F_{\lambda}\) , where \(F_{\lambda}=\sum_{i\in I_{\lambda}}E_{i}\) (``associativity'')
\end{enumerate}

of the ordinal sum). If I is the linearly ordered set \(\{1, 2\}\) , we write \(E_{1} + E_{2}\) for the ordinal sum of \(E_{1}\) and \(E_{2}\) . Show that \(E_{2} + E_{1}\) and \(E_{1} + E_{2}\) are not necessarily isomorphic.

\begin{enumerate}
\def\labelenumi{(\alph{enumi})}
\setcounter{enumi}{2}
\item
  An ordinal sum \(\sum_{i\in I}E_i\) is right directed if and only if I is right directed and \(E_{\omega}\) is right directed for each maximal element \(\omega\) of I.
\item
  An ordinal sum \(\sum_{i\in I}E_i\) is totally ordered if and only if \(I\) and each \(E_i\) is totally ordered.
\item
  An ordinal sum \(\sum_{i\in I}E_i\) is a lattice if and only if the following conditions are satisfied:
\end{enumerate}

\begin{enumerate}
\def\labelenumi{(\Roman{enumi})}
\item
  The set I is a lattice, and for each pair \((\lambda, \mu)\) of non-comparable indices in I, \(E_{\text{sup}(\lambda, \mu)}\) (resp. \(E_{\text{inf}(\lambda, \mu)}\) ) has a least (resp. greatest) element.
\item
  For each \(\alpha \in I\) and each pair \((x, y)\) of elements of \(E_{\alpha}\) such that the set \(\{x, y\}\) is bounded above (resp. bounded below) in \(E_{\alpha}\) , the set \(\{x, y\}\) has a least upper bound (resp. greatest lower bound) in \(E_{\alpha}\) .
\item
  For each \(\alpha \in I\) such that \(E_{\alpha}\) contains a set of two elements which has no upper bound (resp. lower bound) in \(E_{\alpha}\) , the set of indices \(\lambda \in I\) such that \(\lambda > \alpha\) (resp. \(\lambda < \alpha\) ) has a least element (resp. greatest element) \(\beta\) , and \(E_{\beta}\) has a least element (resp. greatest element).
\end{enumerate}

* 4. Let E be an ordered set, and let \((\mathbf{E}_t)_{t \in I}\) be the partition of E formed by the connected components of E (Chapter II, § 6, Exercise 10) with respect to the reflexive and symmetric relation ``either \(x = y\) or \(x\) and \(y\) are not comparable''.

\begin{enumerate}
\def\labelenumi{(\alph{enumi})}
\item
  Show that if \(\iota \neq \kappa\) and \(x \in E_{\iota}\) and \(y \in E_{\kappa}\) , then \(x, y\) are comparable; and that if, for example, \(x \leqslant y\) , \(y' \in E_{\kappa}\) , and \(y' \neq y\) , then also \(x \leqslant y'\) (use the fact that there exists no partition of \(E_{\kappa}\) into two sets A and B such that every element of A is comparable with every element of B).
\item
  Deduce from (a) that the equivalence relation S corresponding to the partition \((\mathbf{E}_{t})\) of E is compatible (in x and in y) with the order relation \(x \leqslant y\) on E, and that the quotient ordered set E/S (Exercise 2) is totally ordered.
\item
  What are the connected components of an ordered set \(\mathbf{E} = \mathbf{F} \times \mathbf{G}\) which is the product of two totally ordered sets? *
\end{enumerate}

\begin{enumerate}
\def\labelenumi{\arabic{enumi}.}
\setcounter{enumi}{4}
\item
  Let E be an ordered set. A subset X of E is said to be free if no two distinct elements of X are comparable. Let \(\Im\) be the set of free subsets of E. Show that, on \(\Im\) , the relation ``given any \(x \in X\) , there exists \(y \in Y\) such that \(x \leqslant y\) is an order relation between X and Y, written \(X \leqslant Y\) . The mapping \(x \to \{x\}\) is an isomorphism of E onto a subset of the ordered set \(\mathfrak{J}\) . If \(X \subset Y\) , where \(X \in \mathfrak{J}\) and \(Y \in \mathfrak{J}\) , show that \(X \leqslant Y\) . The ordered set \(\mathfrak{J}\) is totally ordered if and only if \(E\) is totally ordered, and then \(\mathfrak{J}\) is canonically isomorphic to \(E\) .
\item
  Let E and F be two ordered sets, and let \(\mathcal{A}(\mathbf{E},\mathbf{F})\) be the subset of the product ordered set \(F^{E}\) consisting of the increasing mappings of E into F.
\end{enumerate}

\begin{enumerate}
\def\labelenumi{(\alph{enumi})}
\item
  Show that if E, F, G are there ordered sets, then the ordered set \(\mathcal{A}(\mathbf{E}, \mathbf{F} \times \mathbf{G})\) is isomorphic to the product ordered set \(\mathcal{A}(\mathbf{E}, \mathbf{F}) \times \mathcal{A}(\mathbf{E}, \mathbf{G})\) .
\item
  Show that if E, F, G are three ordered sets, then the ordered set \(\mathcal{A}(\mathbf{E} \times \mathbf{F}, \mathbf{G})\) is isomorphic to the ordered set \(\mathcal{A}(\mathbf{E}, \mathcal{A}(\mathbf{F}, \mathbf{G}))\) .
\item
  If \(\mathbf{E} \neq \emptyset\) , then \(\mathcal{A}(\mathbf{E}, \mathbf{F})\) is a lattice if and only if \(\mathbf{F}\) is a lattice.
\item
  Suppose that E and F are both non-empty. Then \(\mathcal{A}(\mathbf{E}, \mathbf{F})\) is totally ordered if and only if one of the following conditions is satisfied:
\end{enumerate}

(α) F consists of a single element;

(β) E consists of a single element and F is totally ordered;

(γ) E and F are both totally ordered and F has two elements.

\begin{enumerate}
\def\labelenumi{\arabic{enumi}.}
\setcounter{enumi}{6}
\item
  In order that every mapping of an ordered set E into an ordered set F with at least two elements, which is both an increasing and a decreasing mapping, should be constant on E, it is necessary and sufficient that E should be connected with respect to the reflexive and symmetric relation ``x and y are comparable'' (Chapter II, § 6, Exercise 10). This condition is satisfied in particular if E is either left or right directed.
\item
  Let E and F be two ordered sets, let f be an increasing mapping of E into F, and g an increasing mapping of F into E. Let A (resp. B) be the set of all \(x \in E\) (resp. \(y \in F\) ) such that \(g(f(x)) = x\) (resp. \(f(g(y)) = y\) ). Show that the two ordered sets A and B are canonically isomorphic.
\end{enumerate}

* 9. If E is a lattice, prove that

\[
\sup _ {j} \left(\inf _ {i} x _ {i j}\right) \leqslant \inf _ {i} \left(\sup _ {j} x _ {i j}\right)
\]

for every finite ``double'' family \((x_{ij})\) . *

\begin{enumerate}
\def\labelenumi{\arabic{enumi}.}
\setcounter{enumi}{9}
\tightlist
\item
  Let E and F be two lattices. Then a mapping f of E into F is increasing if and only if
\end{enumerate}

\[
f (\inf (x, y)) \leqslant \inf (f (x), f (y))
\]

for all \(x \in \mathbf{E}\) and all \(y \in \mathbf{E}\) .

* Give an example of an increasing mapping \(f\) of the product ordered set \(\mathbf{N} \times \mathbf{N}\) into the ordered set \(\mathbf{N}\) such that the relation

\[
f (\inf (x, y)) = \inf (f (x, y))
\]

is false for at least one pair \((x, y) \in \mathbf{N} \times \mathbf{N}_{*}\) .

\begin{enumerate}
\def\labelenumi{\arabic{enumi}.}
\setcounter{enumi}{10}
\tightlist
\item
  A lattice E is said to be complete if every subset of E has a least upper bound and a greatest lower bound in E; this means, in particular, that E has a greatest and a least element.
\end{enumerate}

\begin{enumerate}
\def\labelenumi{(\alph{enumi})}
\item
  Show that if an ordered set \(\mathbf{E}\) is such that every subset of \(\mathbf{E}\) has a least upper bound in \(\mathbf{E}\) , then \(\mathbf{E}\) is a complete lattice.
\item
  A product of ordered sets is a complete lattice if and only if each of the factors is a complete lattice.
\item
  An ordinal sum (Exercise 3) \(\sum_{i\in I}E_i\) is a complete lattice if and only if the following conditions are satisfied:
\end{enumerate}

\begin{enumerate}
\def\labelenumi{(\Roman{enumi})}
\item
  I is a complete lattice.
\item
  If \(J\) is a subset of \(I\) which has no greatest element, and if \(\sigma = \sup J\) , then \(E_{\sigma}\) has a least element.
\item
  For each \(t \in I\) every subset of \(E_t\) which has an upper bound in \(E_t\) has a least upper bound in \(E_t\) .
\item
  For each \(\iota\in I\) such that \(E_{\iota}\) has no greatest element, the set of all \(x>\iota\) has a least element \(\alpha\) , and \(E_{\alpha}\) has a least element.
\end{enumerate}

\begin{enumerate}
\def\labelenumi{(\alph{enumi})}
\setcounter{enumi}{3}
\tightlist
\item
  The ordered set \(\mathcal{A}(\mathbf{E},\mathbf{F})\) of increasing mappings of an ordered set E into an ordered set F (Exercise 6) is a complete lattice if and only if F is a complete lattice.
\end{enumerate}

\begin{enumerate}
\def\labelenumi{\arabic{enumi}.}
\setcounter{enumi}{11}
\item
  Let \(\Phi\) be a set of mappings of a set A into itself. Let \(\mathfrak{F}\) be the subset of \(\mathfrak{P}(A)\) consisting of all \(X \subset A\) such that \(f(X) \subset X\) for each \(f \in \Phi\) . Show that \(\mathfrak{F}\) is a complete lattice with respect to the relation of inclusion.
\item
  Let E be an ordered set. A mapping f of E into itself is said to be a closure if it satisfies the following conditions: (1) f is increasing; (2) for each \(x \in E\) , \(f(x) \geqslant x\) ; (3) for each \(x \in E\) , \(f(f(x)) = f(x)\) . Let F be the set of elements of E which are invariant under f.
\end{enumerate}

\begin{enumerate}
\def\labelenumi{(\alph{enumi})}
\item
  Show that for each \(x \in E\) the set \(F_{x}\) of elements \(y \in F\) such that \(x \leqslant y\) is not empty and has a least element, namely \(f(x)\) . Conversely, if G is a subset of E such that, for each \(x \in E\) , the set of all \(y \in G\) such that \(x \leqslant y\) has a least element \(g(x)\) , then g is a closure and G is the set of elements of E which are invariant under g.
\item
  Suppose that E is a complete lattice. Show that the greatest lower bound in E of any non-empty subset of F belongs to F.
\item
  Show that if E is a lattice, then \(f(\sup(x,y)) = f(\sup(f(x),y))\) for each pair of elements x, y of E.
\end{enumerate}

\begin{enumerate}
\def\labelenumi{\arabic{enumi}.}
\setcounter{enumi}{13}
\item
  Let A and B be two sets, and let R be any subset of \(A \times B\) . For each subset X of A (resp. each subset Y of B) let \(\rho(X)\) (resp. \(\sigma(Y)\) ) denote the set of all \(y \in B\) (resp. \(x \in A\) ) such that \((x, y) \in R\) for all \(x \in A\) (resp. \((x, y) \in R\) for all \(y \in B\) ). Show that \(\rho\) and \(\sigma\) are decreasing mappings, and that the mappings \(X \to \sigma(\rho(X))\) and \(Y \to \rho(\sigma(Y))\) are closures (Exercise 13) in \(\mathfrak{P}(A)\) and \(\mathfrak{P}(B)\) respectively (ordered by inclusion).
\item
  \begin{enumerate}
  \def\labelenumii{(\alph{enumii})}
  \tightlist
  \item
    Let E be an ordered set, and for each subset X of E let \(\rho(X)\) (resp. \(\sigma(X)\) ) denote the set of upper (resp. lower) bounds of X in E. Show that, in \(\mathfrak{P}(E)\) , the set \(\tilde{E}\) of subsets X such that \(X = \sigma(\rho(X))\) is a complete lattice, and that the mapping \(i: x \to \sigma(\{x\})\) is an isomorphism (called canonical) of E onto an ordered subset \(E'\) of \(\tilde{E}\) such that, if a family ( \(x_t\) ) of elements of E has a least upper bound (resp. greatest lower bound) in E, the image of this least upper bound (resp. greatest lower bound) is the least upper bound (resp. greatest lower bound) in \(\tilde{E}\) of the family of images of the \(x_t\) . \(\tilde{E}\) is called the completion of the ordered set E.
  \end{enumerate}
\end{enumerate}

\begin{enumerate}
\def\labelenumi{(\alph{enumi})}
\setcounter{enumi}{1}
\item
  Show that, for every subset X of E, \(\sigma(\rho(X))\) is the least upper bound in \(\tilde{E}\) of the subset \(i(X)\) of \(\tilde{E}\) . If f is any increasing mapping of E into a complete lattice F, there exists a unique increasing mapping \(\bar{f}\) of \(\tilde{E}\) into F such that \(f = \bar{f} \circ i\) and \(\bar{f}(\sup Z) = \sup (\bar{f}(Z))\) for every subset Z of \(\tilde{E}\) .
\item
  If E is totally ordered, show that \(\tilde{E}\) is totally ordered.
\end{enumerate}

¶ 16. A lattice E is said to be distributive if it satisfies the following two conditions:

\[
\begin{array}{l l} \left(\mathbf {D} ^ {\prime}\right) & \sup (x, \inf (y, z)) = \inf (\sup (x, y), \sup (x, z)), \\ \left(\mathbf {D} ^ {\prime \prime}\right) & \inf (x, \sup (y, z)) = \sup (\inf (x, y) \inf (x, z)) \end{array}
\]

for all \(x, y, z\) in E. A totally ordered set is a distributive lattice.

\begin{enumerate}
\def\labelenumi{(\alph{enumi})}
\tightlist
\item
  Show that each of the conditions (D), (D') separately implies the condition
\end{enumerate}

\[
\begin{array}{l} \text {(D)} \sup (\inf (x, y), \inf (y, z), \inf (z, x)) \\ = \inf (\sup (x, y), \sup (y, z), \sup (z, x)) \\ \text {for all x,y,z in E.} \end{array}
\]

\begin{enumerate}
\def\labelenumi{(\alph{enumi})}
\setcounter{enumi}{1}
\tightlist
\item
  Show that the condition (D) implies the condition
\end{enumerate}

\[
\text {   If   } x \geqslant z, \text {   then   } \sup (z, \inf (x, y)) = \inf (x, \sup (y, z)).\tag{M}
\]

Deduce that (D) implies each of (D') and (D'\,`), and hence that the three axioms (D), (D'), and (D'\,`) are equivalent (to show, for example, that (D) implies (D'), take the least upper bound of x and each side of (D), and use (M)).

\begin{enumerate}
\def\labelenumi{(\alph{enumi})}
\setcounter{enumi}{2}
\tightlist
\item
  Show that each of the two conditions
\end{enumerate}

\[
\begin{array}{r l} \inf (z, \sup (x, y)) & \leqslant \sup (x, \inf (y, z)), \\ \inf (\sup (x, y), \sup (z, \inf (x, y))) & = \sup (\inf (x, y), \inf (y, z), \inf (z, x)) \end{array}
\]

for all \(x, y, z\) in \(\mathbf{E}\) is necessary and sufficient for \(\mathbf{E}\) to be distributive. (To show that \((\mathbf{T}')\) implies \((\mathbf{D}^n)\) , consider the element

\[
\inf (z, \sup (x, \inf (y, z))).)
\]

\begin{enumerate}
\def\labelenumi{\arabic{enumi}.}
\setcounter{enumi}{16}
\tightlist
\item
  A lattice E which has a least element \(\alpha\) is said to be relatively complemented if, for each pair of elements \(x, y\) of E such that \(x \leqslant y\) , there exists an element \(x'\) such that \(\sup (x, x') = y\) and \(\inf (x, x') = \alpha\) . Such an element \(x'\) is called a relative complement of \(x\) with respect to \(y\) .
\end{enumerate}

* (a) Show that the set E of vector subspaces of a vector space of dimension ≥ 2, ordered by inclusion, is a relatively complemented lattice, but that if x, y are two elements of E such that x ≤ y, there exist in general several distinct relative complements of x with respect to y. *

\begin{enumerate}
\def\labelenumi{(\alph{enumi})}
\setcounter{enumi}{1}
\item
  If E is distributive and relatively complemented, show that if \(x \leqslant y\) in E, there exists a unique relative complement of x with respect to y. E is said to be a Boolean lattice if it is distributive and relatively complemented and if, moreover, it has a greatest element \(\omega\) . For each \(x \in E\) , let \(x^{*}\) be the complement of x with respect to \(\omega\) . Then the mapping \(x \to x^{*}\) is an isomorphism of E onto the ordered set obtained by endowing E with the opposite ordering, and we have \((x^{*})^{*} = x\) . If A is any set, then the set \(\mathfrak{P}(A)\) of all subsets of A, ordered by inclusion, is a Boolean lattice.
\item
  If E is a complete Boolean lattice (Exercise 11), show that for each family \((x_{\lambda})\) of elements of E and each \(y \in E\) we have
\end{enumerate}

\[
\inf (y, \sup _ {\lambda} (x _ {\lambda})) = \sup _ {\lambda} (\inf (y, x _ {\lambda})).
\]

(Reduce to the case \(y = \alpha\) , and use the fact that if \(\inf (z, x_{\lambda}) = \alpha\) for every index \(\lambda\) , then \(z^{*} \geqslant x_{\lambda}\) for every \(\lambda\) ).

¶ * 18. Let A be a set with at least three elements, let ℜ be the set of all partitions of A, ordered by the relation ``ω is finer than ω'\,'' between ω and ω' (no. 1, Example 4). Show that ℜ is a complete lattice (Exercise 11), is not distributive (Exercise 17), but is relatively complemented. (To prove the last assertion, well-order the sets belonging to a partition.) *

\begin{enumerate}
\def\labelenumi{\arabic{enumi}.}
\setcounter{enumi}{18}
\tightlist
\item
  An ordered set E is said to be without gaps if it contains two distinct comparable elements and if, for each pair of elements x, y such that x \textless{} y, the open interval {]}x, y{[} is not empty. Show that an ordinal sum \(\sum_{i\in I}E_{i}\) (Exercise 3) is without gaps if and only if the following conditions are satisfied :
\end{enumerate}

\begin{enumerate}
\def\labelenumi{(\Roman{enumi})}
\item
  Either I contains two distinct comparable elements, or else there exists \(\iota \in I\) such that \(E_{t}\) contains two distinct comparable elements.
\item
  Each \(E_{t}\) which contains at least two distinct comparable elements is without gaps.
\item
  If \(\alpha, \beta\) are two elements of \(I\) such that \(\alpha < \beta\) , and if the interval \([]\alpha, \beta[\) in \(I\) is empty, then either \(E_{\alpha}\) has no maximal element or else \(E_{\beta}\) has no minimal element.
\end{enumerate}

In particular, every ordinal sum \(\sum_{i\in I}E_i\) of sets without gaps is itself without gaps provided that no \(E_i\) has a maximal element (or provided that no \(E_i\) has a minimal element). If I is without gaps, and if each \(E_i\) is either without gaps or contains no two distinct comparable elements, then \(\sum E_i\) is without gaps.

\begin{enumerate}
\def\labelenumi{\arabic{enumi}.}
\setcounter{enumi}{19}
\tightlist
\item
  An ordered set E is said to be scattered if no ordered subset of E is without gaps (Exercise 19). Every subset of a scattered set is scattered. * Every well-ordered set of more than one element is scattered. *
\end{enumerate}

\begin{enumerate}
\def\labelenumi{(\alph{enumi})}
\item
  Suppose that E is scattered. Then, if x, y are any two elements of E such that x \textless{} y, there exist two elements \(x'\) , \(y'\) of E such that \(x \leqslant x' < y' \leqslant y\) and such that the interval \(]x', y'[\) is empty. * Give an example of a totally ordered set which satisfies this condition and is not scattered (consider Cantor's triadic set). *
\item
  An ordinal sum \(E = \sum_{i \in I} E_i\) (where neither I nor any \(E_i\) is empty) is scattered if and only if I and each \(E_i\) is scattered. (Note that E contains a subset isomorphic to I and that every subset F of E is the ordinal sum of those sets \(F \cap E_i\) which are non-empty; finally use Exercise 19.)
\end{enumerate}

\begin{enumerate}
\def\labelenumi{\arabic{enumi}.}
\setcounter{enumi}{20}
\item
  Let E be a non-empty totally ordered set, and let \(S\{x,y\}\) be the relation ``the closed interval with endpoints x, y is scattered'' (Exercise 20). Show that S is an equivalence relation which is weakly compatible (Exercise 2) in x and y with the order relation on E, that the equivalence classes with respect to S are scattered sets, and that the quotient ordered set E/S (Exercise 2) is either without gaps or else consists of a single element. Deduce that E is isomorphic to an ordinal sum of scattered sets whose index set is either without gaps or else consists of a single element.
\item
  \begin{enumerate}
  \def\labelenumii{(\alph{enumii})}
  \tightlist
  \item
    Let E be an ordered set. A subset U of E is said to be open if, for each \(x \in U\) , U contains the interval \([x, \rightarrow[\) . An open set U is said to be regular if there exists no open set \(V \supset U\) , distinct from U, such that U is cofinal in V. Show that every open set U is cofinal in exactly one regular open set \(\tilde{U} (*)\) . The mapping \(U \to \tilde{U}\) is increasing. If U, V are two open sets such that \(U \cap V = \varnothing\) , then also \(\tilde{U} \cap \tilde{V} = \varnothing\) .
  \end{enumerate}
\end{enumerate}

\begin{enumerate}
\def\labelenumi{(\alph{enumi})}
\setcounter{enumi}{1}
\item
  Show that the set R(E) of regular open subsets of E, ordered by inclusion, is a complete Boolean lattice (Exercise 17). For R(E) to consist of two elements it is necessary and sufficient that E should be non-empty and right directed.
\item
  If \(\mathbf{F}\) is a cofinal subset of \(\mathbf{E}\) , show that the mapping \(\mathbf{U} \to \mathbf{U} \cap \mathbf{F}\) is an isomorphism of \(\mathbf{R}(\mathbf{E})\) onto \(\mathbf{R}(\mathbf{F})\) .
\item
  If \(E_{1}\) , \(E_{2}\) are two ordered sets, then every open set in \(E_{1} \times E_{2}\) is of the form \(U_{1} \times U_{2}\) , where \(U_{i}\) is open in \(E_{i}\) (i=1, 2). The set \(\mathbf{R}(\mathbf{E}_{1} \times \mathbf{E}_{2})\) is isomorphic to \(\mathbf{R}(\mathbf{E}_{1}) \times \mathbf{R}(\mathbf{E}_{2})\) .
\end{enumerate}

\begin{enumerate}
\def\labelenumi{\arabic{enumi}.}
\setcounter{enumi}{22}
\tightlist
\item
  Let E be an ordered set and let \(\mathbf{R}_{0}(\mathbf{E}) = \mathbf{R}(\mathbf{E}) - \{\phi\}\) (Exercise 22). For each \(x \in E\) , let \(r(x)\) denote the unique regular open set in which the interval \([x, \rightarrow[\) (which is an open set) is cofinal. The mapping r so defined is called the canonical mapping of E into \(\mathbf{R}_{0}(\mathbf{E})\) . Endow \(\mathbf{R}_{0}(\mathbf{E})\) with the order relation opposite to the relation of inclusion.
\end{enumerate}

\begin{enumerate}
\def\labelenumi{(\alph{enumi})}
\item
  Show that the mapping \(r\) is increasing and that \(r(\mathbf{E})\) is cofinal in \(\mathbf{R}_0(\mathbf{E})\) .
\item
  An ordered set E is said to be antidirected if the canonical mapping \(r: \mathbf{E} \to \mathbb{R}_{0}(\mathbf{E})\) is injective. For this to be so it is necessary and sufficient that the following two conditions should be satisfied:
\end{enumerate}

\begin{enumerate}
\def\labelenumi{(\Roman{enumi})}
\item
  If \(x\) and \(y\) are two elements of \(\mathbf{E}\) such that \(x < y\) , there exists \(z \in \mathbf{E}\) such that \(x < z\) and such that the intervals \([y, \rightarrow[\) and \([z, \rightarrow[\) do not intersect.
\item
  If \(x\) and \(y\) are two non-comparable elements of \(\mathbf{E}\) , then either there exists \(x' \geqslant x\) such that the intervals \([x', \rightarrow[\) and \([y, \rightarrow[\) do not intersect, or else there exists \(y' \geqslant y\) such that the intervals \([x, \rightarrow[\) and \([y', \rightarrow[\) do not intersect.
\end{enumerate}

\begin{enumerate}
\def\labelenumi{(\alph{enumi})}
\setcounter{enumi}{2}
\tightlist
\item
  Show that, for every ordered set E, \(R_{0}(E)\) is antidirected and that the canonical mapping of \(R_{0}(E)\) into \(R_{0}(R_{0}(E))\) is bijective (use Exercise 22 (a)).
\end{enumerate}

* 24. (a) An ordered set E is said to be branched (on the right) if for each \(x \in \mathbf{E}\) there exist \(y, z\) in E such that \(x \leqslant y, x \leqslant z\) and the

intervals \([y, \rightarrow[\) and \([z, \rightarrow[\) do not intersect. An antidirected set with no maximal elements (Exercise 23) is branched.

\begin{enumerate}
\def\labelenumi{(\alph{enumi})}
\setcounter{enumi}{1}
\tightlist
\item
  Let E be the set of intervals in R of the form
\end{enumerate}

\[
[ k. 2 ^ {- n}, (k + 1) 2 ^ {- n} ] (0 \leqslant k <   2 ^ {n}),
\]

ordered by the relation \(\supset\) . Show that E is antidirected and has no maximal elements.

\begin{enumerate}
\def\labelenumi{(\alph{enumi})}
\setcounter{enumi}{2}
\item
  Give an example of a branched set in which there exists no anti-directed cofinal subset. (Take the product of the set E defined in (b) with a well-ordered set which contains no countable cofinal subset, and use Exercise 22.)
\item
  Give an example of an ordered set E which is not antidirected, but which has an antidirected cofinal subset (note that an ordinal sum \(\sum_{\xi\in E} F_{\xi}\) contains a cofinal subset isomorphic to E). *
\end{enumerate}

